% Options for packages loaded elsewhere
\PassOptionsToPackage{unicode}{hyperref}
\PassOptionsToPackage{hyphens}{url}
%
\documentclass[
]{article}
\usepackage{amsmath,amssymb}
\usepackage{iftex}
\ifPDFTeX
  \usepackage[T1]{fontenc}
  \usepackage[utf8]{inputenc}
  \usepackage{textcomp} % provide euro and other symbols
\else % if luatex or xetex
  \usepackage{unicode-math} % this also loads fontspec
  \defaultfontfeatures{Scale=MatchLowercase}
  \defaultfontfeatures[\rmfamily]{Ligatures=TeX,Scale=1}
\fi
\usepackage{lmodern}
\ifPDFTeX\else
  % xetex/luatex font selection
\fi
% Use upquote if available, for straight quotes in verbatim environments
\IfFileExists{upquote.sty}{\usepackage{upquote}}{}
\IfFileExists{microtype.sty}{% use microtype if available
  \usepackage[]{microtype}
  \UseMicrotypeSet[protrusion]{basicmath} % disable protrusion for tt fonts
}{}
\makeatletter
\@ifundefined{KOMAClassName}{% if non-KOMA class
  \IfFileExists{parskip.sty}{%
    \usepackage{parskip}
  }{% else
    \setlength{\parindent}{0pt}
    \setlength{\parskip}{6pt plus 2pt minus 1pt}}
}{% if KOMA class
  \KOMAoptions{parskip=half}}
\makeatother
\usepackage{xcolor}
\usepackage[margin=1in]{geometry}
\usepackage{graphicx}
\makeatletter
\def\maxwidth{\ifdim\Gin@nat@width>\linewidth\linewidth\else\Gin@nat@width\fi}
\def\maxheight{\ifdim\Gin@nat@height>\textheight\textheight\else\Gin@nat@height\fi}
\makeatother
% Scale images if necessary, so that they will not overflow the page
% margins by default, and it is still possible to overwrite the defaults
% using explicit options in \includegraphics[width, height, ...]{}
\setkeys{Gin}{width=\maxwidth,height=\maxheight,keepaspectratio}
% Set default figure placement to htbp
\makeatletter
\def\fps@figure{htbp}
\makeatother
\setlength{\emergencystretch}{3em} % prevent overfull lines
\providecommand{\tightlist}{%
  \setlength{\itemsep}{0pt}\setlength{\parskip}{0pt}}
\setcounter{secnumdepth}{-\maxdimen} % remove section numbering
\usepackage{float}
\ifLuaTeX
  \usepackage{selnolig}  % disable illegal ligatures
\fi
\usepackage{bookmark}
\IfFileExists{xurl.sty}{\usepackage{xurl}}{} % add URL line breaks if available
\urlstyle{same}
\hypersetup{
  pdftitle={Homework 1 Solutions},
  hidelinks,
  pdfcreator={LaTeX via pandoc}}

\title{Homework 1 Solutions}
\author{}
\date{\vspace{-2.5em}}

\begin{document}
\maketitle

\({\bf (1)}\) Define the following terms:

\begin{itemize}
\tightlist
\item
  \textbf{(a)} Population - the collection of all possible persons,
  events, or objects of interest
\item
  \textbf{(b)} Sample - the subset of the population that we actually
  observe (the observed observations)
\item
  \textbf{(c)} Parameter - A numerical characteristic of a population
  (alternatively: A function of the entire population)
\item
  \textbf{(d)} Statistic - A numerical characteristic of a sample
  (alternatively: A function of the sample)
\end{itemize}

\({\bf (2)}\) Explain the difference between a \textbf{qualitative}
(categorical) and a \textbf{quantitative} variable

{A qualitative variable is a non-numeric characteristic such as a name
or label while a quantitative variable is numerical characteristic such
as a measurement or count}

\({\bf (3)}\) Explain the difference between a \textbf{discrete} and a
\textbf{continuous} variable and give an example of each

{A discrete variable is a quantitative variable that takes on only
distinct whole number (integer) values such as counts whereas a
continuous variable is a quantitative variable that can assume any value
within a certain interval such temperature or height}

\({\bf (4)}\) Explain the difference between a \textbf{nominal} and an
\textbf{ordinal} variable and give an example of each

{A nominal variable is a qualitative variable that has categories with
no inherent ordering such as the names of a city or the brands or shoes.
An ordinal variable is a qualitative variable that has an inherent
ordering to its categories such as a Likert score or education level.}

\({\bf (5)}\) \textbf{At What age did women marry?} A historian wants to
estimate the average age at marriage of women in New England in the
early 19th century. Within her state archives, she finds marriage
records for the years \(1800 - 1820\), which she treats as a sample of
all marriage records from the early 19th century. The average age of the
women in the records is \(24.1\) years of age. Using the appropriate
statistical method, she estimates that the average age of brides in the
early \(19^{th}\)- century New England was between \(23.5\) and \(24.7\)
years of age.

\begin{itemize}
\tightlist
\item
  \textbf{(a)} Which part of this example gives a descriptive summary of
  the data? {When she computes a mean of \(24.1\) years of age from the
  marriage records}
\item
  \textbf{(b)} Which part of this example draws an inference about a
  population? {The part where she estimates the average age of women to
  be married in the 19th century to be between \(23.5\) and \(24.7\)
  years for the population}
\item
  \textbf{(c)} What population is the historian studying? {All women who
  married in New England in the early 19th century}
\item
  \textbf{(d)} The average age the historian computed from the
  historical records was 24.1 years of age. Is 24.1 years of age a
  statistic or a parameter? Why? {A statistic. The quantity \(24.1\) was
  computed from a sample of marriage records spanning the first \(20\)
  years of the century}
\end{itemize}

\({\bf (6)}\) Consider the following data obtained from flipping a coin
15 times. A value of \(H\) denotes the coin was heads whereas a value of
\(T\) denotes the coin flip was tails.

\begin{table}[H]
\centering\centering
\begin{tabular}{llllllllllllllll}
\toprule
\cellcolor{gray!10}{coin.flip.result.} & \cellcolor{gray!10}{H} & \cellcolor{gray!10}{H} & \cellcolor{gray!10}{H} & \cellcolor{gray!10}{T} & \cellcolor{gray!10}{H} & \cellcolor{gray!10}{T} & \cellcolor{gray!10}{T} & \cellcolor{gray!10}{T} & \cellcolor{gray!10}{H} & \cellcolor{gray!10}{H} & \cellcolor{gray!10}{T} & \cellcolor{gray!10}{T} & \cellcolor{gray!10}{T} & \cellcolor{gray!10}{H} & \cellcolor{gray!10}{T}\\
\bottomrule
\end{tabular}
\end{table}

Compute the frequency table for the variable ``coin flip result'' and
answer the following questions:

\begin{table}[H]
\centering\centering
\begin{tabular}{lrl}
\toprule
Result & Frequency & Relative Frequency\\
\midrule
\cellcolor{gray!10}{H} & \cellcolor{gray!10}{7} & \cellcolor{gray!10}{7/15 = 0.47}\\
T & 8 & 8/15 = 0.53\\
\bottomrule
\end{tabular}
\end{table}

\begin{itemize}
\tightlist
\item
  \textbf{(a)} What type of variable is ``Coin Flip Result''?
  {Qualitative nominal variable}
\item
  \textbf{(b)} What proportion of the coin flips were heads? {\(0.47\)
  or \(47\%\)}
\item
  \textbf{(c)} Why do we not compute the cumulative relative frequency
  for this variable? {The variable is nominal and therefore has no
  inherent ordering for summing up the relative frequencies}
\item
  \textbf{(d)} What is the best graphical display for this variable and
  why? {Since there are only two categories (heads or tails) a pie chart
  or bar graph would be sufficient to show the distribution}
\end{itemize}

\({\bf (7)}\) A survey about color preferences reported the age
distribution of the people who responded. Below are the results

\begin{table}[H]
\centering\centering
\begin{tabular}{lllllll}
\toprule
 &  &  &  &  &  & \\
\midrule
\cellcolor{gray!10}{Age Group (Years)} & \cellcolor{gray!10}{1-18} & \cellcolor{gray!10}{19-24} & \cellcolor{gray!10}{25-35} & \cellcolor{gray!10}{36-50} & \cellcolor{gray!10}{51-69} & \cellcolor{gray!10}{70 and over}\\
Counts & 10 & 97 & 70 & 36 & 14 & 5\\
\cellcolor{gray!10}{\textcolor{red}{Relative Frequency}} & \cellcolor{gray!10}{\textcolor{red}{0.04}} & \cellcolor{gray!10}{\textcolor{red}{0.42}} & \cellcolor{gray!10}{\textcolor{red}{0.30}} & \cellcolor{gray!10}{\textcolor{red}{0.16}} & \cellcolor{gray!10}{\textcolor{red}{0.06}} & \cellcolor{gray!10}{\textcolor{red}{0.02}}\\
\bottomrule
\end{tabular}
\end{table}

Use this table to answer parts a - d

\begin{itemize}
\item
  \textbf{(a)} Compute the relative frequency for each age group
\item
  \textbf{(b)} Make a bar graph where the heights of the bars are
  relative frequencies
  \includegraphics{week1homework_solns_files/figure-latex/unnamed-chunk-4-1.pdf}
\item
  \textbf{(c)} Describe the distribution {the distribution is slightly
  right skew indicating that most respondents were between 19 and 35
  years of age}
\item
  \textbf{(d)} Explain why your bar graph is not a histogram {A
  histogram has intervals of equal length}
\end{itemize}

\({\bf (8)}\) Email spam is the curse of the internet. The table below
gives a compilation of the most common types of spam

\begin{table}[H]
\centering\centering
\begin{tabular}{lr}
\toprule
Type of Spam & Percentage\\
\midrule
\cellcolor{gray!10}{Adult} & \cellcolor{gray!10}{14.5}\\
Financial & 16.2\\
\cellcolor{gray!10}{Health} & \cellcolor{gray!10}{7.3}\\
Leisure & 7.8\\
\cellcolor{gray!10}{Products} & \cellcolor{gray!10}{21.0}\\
\addlinespace
Scams & 14.2\\
\bottomrule
\end{tabular}
\end{table}

Use the table to answer the questions a and b

\begin{itemize}
\tightlist
\item
  \textbf{(a)} Report the modal spam category {The modal category is
  ``Products'' because it has the highest frequency}
\item
  \textbf{(b)} Construct a \textbf{Pareto chart} using the table above
  \includegraphics{week1homework_solns_files/figure-latex/unnamed-chunk-6-1.pdf}
\end{itemize}

\({\bf (9)}\) A farmer in Idaho is interested in the number of rainy
days in a given year so he records the number of rainy \(R\), cloudy
\(C\) and sunny \(S\) days over two weeks in May. His observations are
\[\{R,R,C,C,C,S,S,C,R,R,C,R,S,S\}\] Use the farmer's observations to
answer the following questions.

\begin{itemize}
\tightlist
\item
  \textbf{(a)} Based on the farmer's observations, what proportion of
  days were rainy? {\[\hat{p}_{\text{rainy}} = \frac{5}{14} = 0.36\] }
\item
  \textbf{(b)} Is the proportion in part \textbf{(a)} a statistic or a
  parameter? why? {It is a statistic because it was computed from a
  sample of 14 observations of the weather}
\item
  \textbf{(c)} What population is the farmer studying? is it a finite or
  infinite population? {The population is all 365 days in a given year.
  It is finite}
\end{itemize}

\({\bf (10)}\) Consider the following data from a survey conducted on
college students in the state of Florida. As part of their research,
surveyors recorded the high school performance (measured in grade point
average - GPA) of 35 college students from across the state. For your
convenience the data have been sorted from least to greatest:

\[2.0, 2.1, 2.3, 2.8, 3.0, 3.0, 3.0, 3.0, 3.0, 3.0,\]
\[3.0, 3.0, 3.0, 3.1, 3.2, 3.3, 3.3, 3.4,  3.4, 3.4,\]
\[3.4, 3.5, 3.5, 3.5, 3.5, 3.6, 3.6, 3.7, 3.7, 3.8,\]
\[3.8, 3.8, 3.8, 4.0, 4.0\]

The following frequency table gives the distribution of the variable
high school GPA (with values rounded to nearest 0.5). Fill in the table
and answer the following questions:

\begin{table}[H]
\centering\centering
\begin{tabular}{rrrr}
\toprule
GPA & Frequency & Relative frequency & Cumulative RF\\
\midrule
\cellcolor{gray!10}{2.0} & \cellcolor{gray!10}{2} & \cellcolor{gray!10}{0.06} & \cellcolor{gray!10}{0.06}\\
2.5 & 1 & 0.03 & 0.09\\
\cellcolor{gray!10}{3.0} & \cellcolor{gray!10}{12} & \cellcolor{gray!10}{0.34} & \cellcolor{gray!10}{0.43}\\
3.5 & 14 & 0.40 & 0.83\\
\cellcolor{gray!10}{4.0} & \cellcolor{gray!10}{6} & \cellcolor{gray!10}{0.17} & \cellcolor{gray!10}{1.00}\\
\bottomrule
\end{tabular}
\end{table}

\begin{itemize}
\tightlist
\item
  \textbf{(a)} What type of variable is GPA? {Quantitative continuous}
\item
  \textbf{(b)} What proportion of college students had a high school GPA
  \(> 3.0\)? (use the rounded frequency table) {From the rounded
  frequency table, \(14 + 6 = 20\) of the 35 students fall above 3.0, so
  the proportion is \(20/35 = 0.57\). (Working from the raw, unrounded
  data instead gives \(22/35 = 0.63\).)}
\item
  \textbf{(c)} What proportion of college students had a high school GPA
  \(< 3.0\)? (use the rounded frequency table) {From the rounded
  frequency table, \(2 + 1 = 3\) of the 35 students fall below 3.0, so
  the proportion is \(3/35 = 0.09\). (Working from the raw, unrounded
  data instead gives \(4/35 = 0.11\).)}
\item
  \textbf{(d)} What is the \textbf{mean} GPA in this sample? {Using the
  raw data,
  \[\bar{x} = \frac{2.0+2.1+...+4.0+4.0}{35} = \frac{114.5}{35} = 3.27 \approx 3.3\]
  Using the rounded frequency table instead gives a very slightly
  different value,
  \[\bar{x} = \frac{2(2.0)+1(2.5)+12(3.0)+14(3.5)+6(4.0)}{35} = \frac{115.5}{35} = 3.30\]}
\item
  \textbf{(e)} What is the \textbf{median} GPA in this sample? {median
  \(= 3.4\)}
\item
  \textbf{(f)} What is the \textbf{mode} of GPA in this sample?
\end{itemize}

{Using the frequency table on the rounded values the mode is the value
with the highest relative frequency \$ = 3.5\$}

{(Using the raw data) the mode is the most frequent value observed in
the data which is \(3.0\) }

\begin{itemize}
\item
  \textbf{(g)} Construct a dot plot for the variable GPA (hint use the
  frequency table)
  \includegraphics{week1homework_solns_files/figure-latex/unnamed-chunk-8-1.pdf}
\item
  \textbf{(h)} Construct a histogram for the variable GPA (hint use the
  frequency table)
  \includegraphics{week1homework_solns_files/figure-latex/unnamed-chunk-9-1.pdf}
\item
  \textbf{(i)} Construct a stem and leaf plot for the variable GPA
\end{itemize}

\begin{verbatim}
## [1] "Using the Rounded GPA scores:"
\end{verbatim}

\begin{verbatim}
## 
##   The decimal point is 1 digit(s) to the left of the |
## 
##   20 | 00
##   22 | 
##   24 | 0
##   26 | 
##   28 | 
##   30 | 000000000000
##   32 | 
##   34 | 00000000000000
##   36 | 
##   38 | 
##   40 | 000000
## 
## NULL
\end{verbatim}

\begin{verbatim}
## [1] "Using the Raw Data:"
\end{verbatim}

\begin{verbatim}
## 
##   The decimal point is 1 digit(s) to the left of the |
## 
##   20 | 0
##   21 | 0
##   22 | 
##   23 | 0
##   24 | 
##   25 | 
##   26 | 
##   27 | 
##   28 | 0
##   29 | 
##   30 | 000000000
##   31 | 0
##   32 | 0
##   33 | 00
##   34 | 0000
##   35 | 0000
##   36 | 00
##   37 | 00
##   38 | 0000
##   39 | 
##   40 | 00
## 
## NULL
\end{verbatim}

\({\bf (11)}\) Which statistic is more resistant to outliers, the mean
or median? Why?

{The median is resistant to outliers because it relies on only the
middle value (or middle two values if \(n\) is even). Therefore, how far
a value is from the center does not influence the median. The mean,
however, is computed using \textbf{ALL} values in a sample distribution.
This causes the mean to be pulled in the direction of extreme values.}

\({\bf (12)}\) Describe the shape of the following distributions and for
each distribution identify if the mean will be larger, smaller or the
same as the median.

\textbf{(a)} {Skew right. The mean will be greater than median because
the center of the data is pulled toward the right side of the
distribution}

\includegraphics{week1homework_solns_files/figure-latex/unnamed-chunk-11-1.pdf}

\textbf{(b)}\\
{Symmetric/Bell shaped. The mean and the median will be equal}

\includegraphics{week1homework_solns_files/figure-latex/unnamed-chunk-12-1.pdf}

\textbf{(c)}\\
{Skew left. The mean will be less than median because the center of the
data is pulled toward the left side of the distribution}

\includegraphics{week1homework_solns_files/figure-latex/unnamed-chunk-13-1.pdf}

\newpage

\({\bf (13)}\) Consider the following four sets of observations of a
quantitative variable \(x\). For your convenience the observations have
been sorted in increasing order. Match datasets \(1-4\) with the correct
histogram (labeled \(A - D\))

\[\text{Dataset 1} = \{0.1, 1.1, 2.6, 2.7, 3.4, 3.4, 4.1, 4.4, 8.8, 9.6\}\]
\[\text{Dataset 2}= \{0.1, 0.3, 1.2, 2.4, 4.4, 4.5, 8.0, 8.9, 9.3, 9.3\}\]
\[\text{Dataset 3} = \{1.1, 3.8, 5.3, 6.0, 6.2, 6.9, 7.9, 7.9, 8.1, 8.7\}\]
\[\text{Dataset 4} = \{3.4, 4.5, 5.4, 5.6, 7.0, 8.5, 8.9, 9.2, 9.7, 9.7\}\]
\includegraphics{week1homework_solns_files/figure-latex/unnamed-chunk-14-1.pdf}

\begin{itemize}
\tightlist
\item
  {Dataset 1 goes with histogram A}
\item
  {Dataset 2 goes with histogram D}
\item
  {Dataset 3 goes with histogram B}
\item
  {Dataset 4 goes with histogram C}
\end{itemize}

\({\bf (14)}\) Consider the following set of 10 observations of a
variable \(X\) sorted from least to greatest:
\[3.3, 3.8, 4.0, 4.8, 4.8, 5.1, 5.2, 5.6, 5.7, 6.9\] Use the data to
answer parts a-b

\begin{itemize}
\tightlist
\item
  \textbf{(a)} Given the data, compute the \(40^{th}\) percentile of
  \(X\) {With \(n = 10\) observations, the \(40^{th}\) percentile falls
  at position \(0.40 \times 10 = 4\). Because this lands exactly on a
  data position, we average the \(4^{th}\) and \(5^{th}\) values - both
  of which are \(4.8\): \[P_{40} = \frac{4.8+4.8}{2} = 4.8\] (Counting
  the \(4^{th}\) value directly gives the same answer here.)}
\item
  \textbf{(b)} Compute Interquartile Range (IQR) range of \(X\)
  {\[IQR = 5.6 - 4.0 = 1.6\]}
\end{itemize}

\({\bf (15)}\) Consider the following \(n = 20\) observations of several
popular cereal brands and answer questions a - g about the variable
``Sugar (g)'':

\begin{table}[H]
\centering\centering\centering
\begin{tabular}{l|r|r|>{\raggedright\arraybackslash}p{18em}}
\hline
Brand & Sodium (mg) & Sugar (g) & Type\\
\hline
\cellcolor{gray!10}{Frosted Mini Wheats} & \cellcolor{gray!10}{0} & \cellcolor{gray!10}{11} & \cellcolor{gray!10}{A}\\
\hline
Raisin Bran & 340 & 18 & A\\
\hline
\cellcolor{gray!10}{All Bran} & \cellcolor{gray!10}{70} & \cellcolor{gray!10}{5} & \cellcolor{gray!10}{A}\\
\hline
Apple Jacks & 140 & 14 & C\\
\hline
\cellcolor{gray!10}{Captain Crunch} & \cellcolor{gray!10}{200} & \cellcolor{gray!10}{12} & \cellcolor{gray!10}{C}\\
\hline
Cheerios & 180 & 1 & C\\
\hline
\cellcolor{gray!10}{Cinnamon Toast Crunch} & \cellcolor{gray!10}{210} & \cellcolor{gray!10}{10} & \cellcolor{gray!10}{C}\\
\hline
Crackling Oat Bran & 150 & 16 & A\\
\hline
\cellcolor{gray!10}{Fiber One} & \cellcolor{gray!10}{100} & \cellcolor{gray!10}{0} & \cellcolor{gray!10}{A}\\
\hline
Frosted Flakes & 130 & 12 & C\\
\hline
\cellcolor{gray!10}{Froot Loops} & \cellcolor{gray!10}{140} & \cellcolor{gray!10}{14} & \cellcolor{gray!10}{C}\\
\hline
Honey Bunches of Oats & 180 & 7 & A\\
\hline
\cellcolor{gray!10}{Honey Nut Cheerios} & \cellcolor{gray!10}{190} & \cellcolor{gray!10}{9} & \cellcolor{gray!10}{C}\\
\hline
Life & 160 & 6 & C\\
\hline
\cellcolor{gray!10}{Rice Krispies} & \cellcolor{gray!10}{290} & \cellcolor{gray!10}{3} & \cellcolor{gray!10}{C}\\
\hline
Honey Smacks & 50 & 15 & A\\
\hline
\cellcolor{gray!10}{Special K} & \cellcolor{gray!10}{220} & \cellcolor{gray!10}{4} & \cellcolor{gray!10}{A}\\
\hline
Wheaties & 180 & 4 & A\\
\hline
\cellcolor{gray!10}{Corn Flakes} & \cellcolor{gray!10}{200} & \cellcolor{gray!10}{3} & \cellcolor{gray!10}{A}\\
\hline
Honeycomb & 210 & 11 & C\\
\hline
\end{tabular}
\end{table}

\begin{itemize}
\tightlist
\item
  \textbf{(a)} Construct a histogram of the data using the following
  bins: \(\leq 2 (g)\), \(2 - 4 (g)\), \(4 - 6 (g)\),
  \(6 - 8 (g)\),\(8 - 10 (g)\), \(10 - 12 (g)\),
  \(12 - 14 (g)\),\(14 - 16 (g)\),\(> 16 (g)\) (hint start by creating a
  frequency table)
\end{itemize}

\begin{figure}
\centering
\includegraphics{week1homework_solns_files/figure-latex/unnamed-chunk-16-1.pdf}
\caption{(A) Histogram of the variable sugar (g) using the bins in the
frequency table (above)}
\end{figure}

\begin{table}[H]
\centering\centering
\caption{\label{tab:unnamed-chunk-16}Frequency table used to construct histogram for the variable "Sugar (g)"}
\centering
\begin{tabular}[t]{lrrr}
\toprule
Sugar (g) & Frequency & RF(x) & CRF(x)\\
\midrule
\cellcolor{gray!10}{< = 2 (g)} & \cellcolor{gray!10}{2} & \cellcolor{gray!10}{0.10} & \cellcolor{gray!10}{0.10}\\
2 - 4 (g) & 4 & 0.20 & 0.30\\
\cellcolor{gray!10}{4 - 6 (g)} & \cellcolor{gray!10}{2} & \cellcolor{gray!10}{0.10} & \cellcolor{gray!10}{0.40}\\
6 - 8 (g) & 1 & 0.05 & 0.45\\
\cellcolor{gray!10}{8 - 10 (g)} & \cellcolor{gray!10}{2} & \cellcolor{gray!10}{0.10} & \cellcolor{gray!10}{0.55}\\
\addlinespace
10 - 12 (g) & 4 & 0.20 & 0.75\\
\cellcolor{gray!10}{12 - 14 (g)} & \cellcolor{gray!10}{2} & \cellcolor{gray!10}{0.10} & \cellcolor{gray!10}{0.85}\\
14 - 16 (g) & 2 & 0.10 & 0.95\\
\cellcolor{gray!10}{> 16 (g)} & \cellcolor{gray!10}{1} & \cellcolor{gray!10}{0.05} & \cellcolor{gray!10}{1.00}\\
\bottomrule
\end{tabular}
\end{table}

\begin{itemize}
\item
  \textbf{(b)} Compute the mean, median, and mode of the variable
  ``Sugar (g)'' {Using the raw data there are 5 modes each with a
  frequency of 2 mode(s) \(= 3,4,11,12,14\)} {median
  \(= \frac{9+10}{2} = 9.5\)}
  {\[\bar{x} = \frac{1}{20}\sum_{i=1}^{20} x_i = \frac{(11+18+5+14+...+3+11)}{20} = 8.75\]
  (note that the mean may be slightly different if you used the
  frequency table above}
\item
  \textbf{(c)} Compute the variance and standard deviation of the
  variable ``Sugar (g)''
  {\[s^2 = \frac{1}{19}\sum_{i=1}^{20} (x_i-\bar{x})^2 = \frac{(11 - 8.75)^2 + (18-8.75)^2 + ...+(3-8.75)^2 + (11 - 8.75)^2}{19}\]
  \[ = \frac{5.06 + 85.56 + ...+33.06+5.06}{19} = 28.3\]
  \[s = \sqrt{s^2} = \sqrt{28.3} = 5.32\]}
\item
  \textbf{(d)} Compute the quartiles (Q1, Q2, Q3) and interquartile
  range of the variable ``Sugar (g)''
  {\[Q1 = 4, \ Q2 = \text{median} = 9.5, \ Q3 = 13, \ IQR = 13 - 4 = 9\]}
\item
  \textbf{(e)} Create a box and whisker plot for the variable ``Sugar
  (g)'' and mark the Q1, Q2, and Q3 quartiles, the median and any
  potential outliers
\end{itemize}

\begin{figure}
\centering
\includegraphics{week1homework_solns_files/figure-latex/unnamed-chunk-17-1.pdf}
\caption{Boxplot of Sugar content in 20 cereal brands}
\end{figure}

\begin{itemize}
\tightlist
\item
  \textbf{(f)} Plot the Cumulative Distribution of the variable ``Sugar
  (g)''
\end{itemize}

\begin{figure}
\centering
\includegraphics{week1homework_solns_files/figure-latex/unnamed-chunk-18-1.pdf}
\caption{Cumulative Distribution of Sugar (g)}
\end{figure}

\begin{itemize}
\tightlist
\item
  \textbf{(g)} What proportion of cereals have a sugar content less than
  or equal to 10 grams? {\[55\%\]}
\end{itemize}

\end{document}
